% !TEX program = xelatex
% twCRPS as the tail arbiter between single-threshold Brier and plain CRPS.
% Table generated by make_spectrum_table.R -> spectrum_table_body.tex.
\documentclass[12pt]{article}

\usepackage{fontspec}
\usepackage{xeCJK}
\setCJKmainfont[AutoFakeSlant=.2, AutoFakeBold=3]{AR PL KaitiM Big5}

\usepackage[margin=1in]{geometry}
\usepackage{amsmath, amssymb}
\usepackage{booktabs}
\usepackage{array}
\usepackage{hyperref}
\hypersetup{colorlinks=true, linkcolor=blue!45!black, urlcolor=blue!45!black}

\title{\textbf{twCRPS 站哪邊？尾部帶裁決\\plain CRPS 的「AR(2) 較差」}}
\author{Brier $\to$ twCRPS $\to$ CRPS 光譜}
\date{2026 年 7 月}

\begin{document}
\maketitle

\section*{問題}

臭氧內插上，AR(2) 相對 AR(1) 出現了一個張力：單門檻 Brier 測不出差異，
但 plain CRPS 顯著較差。plain CRPS 評的是\textbf{整條}預測分布，所以「較差」
可能只是 \emph{bulk}（分布主體）的效應，未必落在超標預測在乎的\textbf{尾部}。
twCRPS 正好能裁決，因為三個分數是\textbf{同一個被積函數的不同讀法}：
\[
  \underbrace{\mathrm{BS}(L)}_{\text{一個門檻}}\;\to\;
  \underbrace{\mathrm{twCRPS}=\!\int_{q_{.90}}^{q_{.995}}\!\mathrm{BS}(z)\,dz}_{\text{尾部帶}}\;\to\;
  \underbrace{\mathrm{CRPS}=\!\int_{-\infty}^{\infty}\!\mathrm{BS}(z)\,dz}_{\text{整條（bulk 主導）}}
\]

\bigskip
% Auto-generated by make_spectrum_table.R -- do not edit by hand.
\begin{table}[htbp]
\centering
\caption{Ozone AR(2) vs.\ AR(1) (skew-$t$, $K{=}7$, $T{=}50$): the same
contrast read by three scores on the Brier$\to$twCRPS$\to$CRPS spectrum,
with $95\%$ site-clustered bootstrap CIs. $\Delta = \mathrm{score}_{AR(2)} -
\mathrm{score}_{AR(1)}$; Brier is dimensionless, twCRPS and CRPS in ppb.
The apparent ``AR(2) worse'' of plain CRPS is a bulk effect: restricting the
integral to the tail band (twCRPS) returns a powerful null, agreeing with
single-threshold Brier.}
\label{tab:twcrps-spectrum}
\small
\begin{tabular}{@{}llccc@{}}
\toprule
Score & Weights & $\Delta$ [95\% CI] & $p$ & reading \\
\midrule
\multicolumn{5}{@{}l}{\emph{200/200 split}}\\
Brier $q_{0.95}$ & one threshold & $-0.000062$ {\scriptsize $[-0.00042,+0.00029]$} & $p{=}0.753$ & ns \\
twCRPS & tail band $[q_{.90},q_{.995}]$ & $-0.00015$ {\scriptsize $[-0.0047,+0.0043]$} & $p{=}0.981$ & ns \\
CRPS & whole line (bulk) & $+0.021$ {\scriptsize $[+0.003,+0.04]$} & $p{=}0.026$ & \textbf{worse} \\
\addlinespace
\multicolumn{5}{@{}l}{\emph{300/100 split}}\\
Brier $q_{0.95}$ & one threshold & $-0.000027$ {\scriptsize $[-0.00033,+0.00032]$} & $p{=}0.808$ & ns \\
twCRPS & tail band $[q_{.90},q_{.995}]$ & $-0.00021$ {\scriptsize $[-0.0044,+0.0047]$} & $p{=}0.868$ & ns \\
CRPS & whole line (bulk) & $+0.039$ {\scriptsize $[+0.0037,+0.072]$} & $p{=}0.027$ & \textbf{worse} \\
\bottomrule
\end{tabular}
\end{table}


\section*{裁決：CRPS 的訊號住在 bulk，不在尾部}

twCRPS 站在 \textbf{Brier 那邊}：兩個 split 的 $\Delta$ 都\textbf{幾乎精確為 $0$}
（$p=0.98$、$0.87$），CI 涵蓋 $0$。而且這不是「沒力氣測」——twCRPS 把整段尾部
的資訊併起來，CI 半寬只有基準 twCRPS（$\approx 0.57$ ppb）的 \textbf{約 $0.8\%$}，
是一個\textbf{有力的虛無}。相對地，plain CRPS 在同樣的資料上顯著較差
（$+0.021$／$+0.039$ ppb，$p\approx0.03$）。

唯一能同時解釋這三格的說法是：\textbf{AR(2) 只稍微弄糟了中央（bulk）預測分布，
plain CRPS 因為被 bulk 主導而偵測到它；但在超標尾部，AR(2) 與 AR(1) 無法區分。}
換句話說，plain CRPS 的「AR(2) 較差」\textbf{對本論文的目標（門檻超標預測）不相關}。

\section*{三個收穫}

\begin{enumerate}
  \item \textbf{twCRPS 不等於單門檻 Brier。}它有\emph{更高的檢定力}（CI 更窄），
        因為把整段尾部的超標資訊併起來，而不是押在單一門檻的少數事件上。
  \item \textbf{twCRPS 也不等於 plain CRPS。}它\emph{對口目標}（尾部加權），能把
        bulk 效應濾掉。這正是它在本例裡與 plain CRPS 分道揚鑣的原因。
  \item \textbf{對論文的結論}：內插任務上，AR(2) 對尾部超標預測\textbf{無益也無害}
        ——這個虛無現在由\emph{兩個}尾部分數（單門檻 Brier 與高檢定力的 twCRPS）
        共同支撐，比單靠 Brier 強得多；而 time-block 外推上 AR(2) 的增益仍然成立。
        「AR(2) 是外推工具、非內插工具」的敘述至此最為乾淨。
\end{enumerate}

\medskip
\noindent\footnotesize 程式：\verb|ozone/US-all-auto/twcrps_bootstrap.R|
（尾部帶 $[q_{.90},q_{.995}]$、$G{=}41$ 網格梯形積分、站層 bootstrap）。
資料：\verb|{brier,twcrps,crps}_bootstrap_pairs_{200-200,300-100}.csv|。

\end{document}
